\chapter{Introduction}
\label{ch:introduction}

\section{Introduction and Clinical Motivation}
Cardiovascular diseases (CVDs) and systemic vascular disorders remain the leading causes of global mortality, necessitating continuous advancements in diagnostic imaging and preventative clinical care. The structural and morphological integrity of the human vascular network serves as a direct biomarker for a multitude of severe pathological conditions, including hypertension, arteriosclerosis, diabetic retinopathy, and various forms of ischemia. In clinical practice, angiography---a specialized medical imaging technique utilized to visualize the inside, or lumen, of blood vessels and organs of the body---has become the gold standard for diagnosing vascular abnormalities. Through the injection of radio-opaque contrast agents and the use of X-ray imaging, fluoroscopy, or advanced optical coherence tomography (OCT), angiographic imaging provides physicians with a high-resolution map of the circulatory system. 

Despite the high fidelity of modern angiographic imaging modalities, the extraction of actionable quantitative metrics from these images remains a significant bottleneck. The precise segmentation and topological extraction of blood vessels from angiography images are of paramount importance. A segmented vascular map allows for the calculation of critical morphological features such as vessel diameter, tortuosity, branching angles, and the total count of distinct vascular networks. These metrics are indispensable for the early detection of microvascular abnormalities, the planning of intricate surgical interventions, such as stent placement and bypass surgeries, and the longitudinal monitoring of disease progression.

Traditionally, the segmentation of blood vessels has been performed manually by highly trained clinicians and radiologists. However, manual annotation is a fundamentally flawed approach in the context of modern, high-throughput clinical environments. It is exceedingly labor-intensive, taking hours to precisely annotate a single high-resolution angiogram. Furthermore, manual segmentation is highly subjective and prone to inter-observer and intra-observer variability. The cognitive fatigue associated with tracing complex, tortuous, and often overlapping vascular structures inevitably leads to human error. Consequently, the development of fully automated, highly precise, and computationally efficient methods for blood vessel segmentation is not merely an academic pursuit but an urgent clinical necessity.

\section{Challenges in Automated Blood Vessel Segmentation}
The transition from manual annotation to automated algorithmic segmentation is fraught with formidable technical challenges. While early attempts utilizing traditional image processing techniques---such as matched filtering, Hessian-based Frangi filters, and morphological operations---showed initial promise, they fundamentally struggle to generalize across diverse clinical datasets. The difficulties in achieving robust automated segmentation stem from several factors inherent to both the biological complexity of the vascular system and the physics of the image acquisition process:

\subsection{High Structural Complexity and Variance}
The human vascular network is an incredibly dense and topologically complex fractal structure. Vessels exhibit extreme variance in scale, ranging from prominent, thick primary arteries and veins to exceptionally thin terminal micro-capillaries that may be only one or two pixels wide in digital representations. Furthermore, vessels branch at varying angles, overlap in two-dimensional projections, and display high degrees of tortuosity. An automated system must possess the representational capacity to simultaneously recognize large, high-contrast structures and preserve the delicate continuity of ultra-thin micro-vessels. Failure to do so results in fragmented segmentations that are useless for downstream topological analysis, such as skeletonization.

\subsection{Severe Class Imbalance}
In the context of binary semantic segmentation, angiographic images present a severe class imbalance problem. The pixels belonging to the foreground class (blood vessels) typically constitute only a small fraction---often less than 10\%---of the total image area, while the vast majority of pixels belong to the background class (tissue, retina, or surrounding anatomical structures). Standard loss functions, such as standard Cross-Entropy, heavily penalize the model based on pixel volume. Consequently, a naive deep learning model trained on such imbalanced data will often collapse into a degenerate state, predicting the entire image as background to achieve an artificially high, yet clinically useless, global accuracy. Specialized loss formulations are absolutely necessary to force the optimization landscape to focus on the sparse foreground structures.

\subsection{Low Contrast and Uneven Illumination}
The acquisition of angiographic images is highly susceptible to physical artifacts. Uneven distribution of the contrast agent, variations in tissue thickness, and the inherent physics of X-ray or optical scattering lead to non-uniform background illumination and localized regions of exceedingly low contrast. Micro-capillaries, in particular, often fade into the background noise, possessing pixel intensities that are statistically indistinguishable from surrounding tissue artifacts. Furthermore, central light reflections and shadowing effects can create false edges that mislead traditional gradient-based edge detection algorithms.

\subsection{Pathological Interferences and Artifacts}
Clinical angiograms are rarely pristine. The very pathologies that necessitate the imaging often introduce severe visual interferences. The presence of lesions, exudates, hemorrhages, and surgical artifacts (such as previously placed stents or clips) introduces significant noise. These pathological structures can visually obscure the underlying vessels or, conversely, mimic their tubular appearance, confusing both traditional algorithms and poorly regularized machine learning models.

\section{The Paradigm Shift to Deep Learning}
To overcome the limitations of traditional, hand-crafted feature engineering, the field of medical image analysis has witnessed a paradigm shift towards Deep Learning (DL), specifically Convolutional Neural Networks (CNNs). Unlike traditional methods that rely on heuristic rules, deep learning models autonomously discover hierarchical, data-driven feature representations directly from annotated clinical datasets. 

The introduction of Fully Convolutional Networks (FCNs) marked the beginning of modern semantic segmentation, allowing networks to output dense, pixel-wise classifications rather than single image-level labels. This was rapidly superseded by the U-Net architecture, which introduced a symmetric encoder-decoder structure with skip connections. These skip connections bypassed the spatial bottleneck, fusing high-resolution, low-level spatial features from the encoder with low-resolution, high-level semantic features from the decoder. U-Net quickly became the de facto standard for medical image segmentation.

However, as clinical demands for precision increased, particularly for the extraction of ultra-thin vessels, the limitations of the standard U-Net became apparent. The semantic gap between the encoder and decoder feature maps often results in the loss of fine-grained spatial details, leading to disconnected vessel segments. This thesis explores advanced architectures designed to bridge this semantic gap and capture multi-scale contextual information more effectively. Specifically, we investigate \textbf{UNet++}, an architecture featuring nested, dense skip pathways, and \textbf{DeepLabV3+}, an architecture utilizing Atrous Spatial Pyramid Pooling (ASPP) to capture features at multiple receptive field scales without degrading spatial resolution.

\section{Research Objectives and Questions}
The overarching goal of this thesis is to engineer a highly accurate, automated, end-to-end pipeline for the semantic segmentation and subsequent morphological analysis of blood vessels from angiography images. To achieve this, the research is guided by the following specific objectives and research questions:

\subsection{Primary Objectives}
\begin{enumerate}
    \item \textbf{Implementation of Advanced Architectures:} To design, implement, and optimize two state-of-the-art deep learning architectures---UNet++ (with and without a pre-trained ResNet50 backbone) and DeepLabV3+---for the specific task of binary vessel segmentation.
    \item \textbf{Mitigation of Class Imbalance:} To formulate, apply, and evaluate advanced hybrid loss functions, specifically combining Binary Cross-Entropy (BCE) with Dice Loss, and Focal Loss with Dice Loss, to counteract the severe foreground-background pixel imbalance inherent in angiograms.
    \item \textbf{Morphological Feature Extraction:} To develop a robust post-processing pipeline utilizing mathematical morphology and skeletonization (thinning) algorithms to extract the centerline topology of the predicted vessel masks, enabling the automated counting of distinct connected vessel networks, branch points, and endpoints.
    \item \textbf{Rigorous Empirical Evaluation:} To systematically evaluate the proposed methodologies on a substantial dataset of 637 annotated angiographic images, comparing the architectures across standard clinical metrics including Intersection over Union (IoU), Dice Coefficient, Precision, and Recall.
\end{enumerate}

\subsection{Key Research Questions}
\begin{itemize}
    \item \textit{RQ1:} How does the dense, nested skip pathway architecture of UNet++ compare to the Atrous Spatial Pyramid Pooling mechanism of DeepLabV3+ in preserving the spatial continuity of ultra-thin micro-capillaries in noisy angiograms?
    \item \textit{RQ2:} To what extent do hybrid loss formulations (BCE+Dice vs. Focal+Dice) accelerate convergence and improve the absolute Intersection over Union (IoU) compared to standard loss functions when dealing with highly imbalanced vascular datasets?
    \item \textit{RQ3:} Can downstream mathematical skeletonization be reliably applied to the probabilistic outputs of deep learning models to extract clinically meaningful, quantitative morphological biomarkers, such as the total count of distinct vessel networks?
\end{itemize}

\section{Thesis Contributions}
This thesis makes several substantial contributions to the fields of medical image analysis and computer vision. By combining state-of-the-art semantic segmentation architectures with rigorous topological post-processing, this work bridges the gap between raw pixel classification and clinically actionable quantitative metrics. The major contributions are enumerated below:

\begin{itemize}
    \item \textbf{Comprehensive Architectural Comparison:} We provide a rigorous, head-to-head empirical evaluation of UNet++ and DeepLabV3+ in the specific domain of angiography. We demonstrate how the nested decoder of UNet++ and the ASPP module of DeepLabV3+ uniquely handle the multi-scale challenges of vessel extraction, providing valuable insights into architectural selection for tubular structure segmentation.
    \item \textbf{Optimized Hybrid Loss Formulations:} We successfully implement and validate hybrid loss functions that force the gradient descent optimization to prioritize overlapping regional predictions (Dice) while simultaneously penalizing hard-to-classify pixels (Focal Loss). We demonstrate that these formulations prevent background collapse and significantly elevate the Dice Coefficient on the validation set.
    \item \textbf{Integration of Heavy Data Augmentation:} We integrate the Albumentations library to apply severe, dynamic spatial and pixel-level augmentations during training. We show that this approach drastically reduces overfitting and improves the generalization of the models on unseen test distributions.
    \item \textbf{Automated Topological Biomarker Extraction:} Moving beyond simple pixel accuracy, we introduce a complete downstream pipeline that utilizes iterative skeletonization algorithms on the predicted masks. This allows the system to autonomously count the number of distinct, disconnected vascular networks and map branch/endpoints, transitioning the output from a visual mask to a quantitative clinical report.
\end{itemize}

\section{Organization of the Thesis}
To provide a clear, logical, and rigorous exposition of the research conducted, the remainder of this thesis is structured as follows:

\textbf{Chapter 2: Literature Review} provides an extensive, chronological review of the state-of-the-art in medical image segmentation. It covers the evolution of algorithms from early hand-crafted filters to the dominance of Convolutional Neural Networks, detailing the architectural innovations of the U-Net family and the DeepLab series. It also reviews the literature on loss functions for imbalanced data and post-processing skeletonization techniques.

\textbf{Chapter 3: Preliminaries and System Model} lays the mathematical and theoretical foundation for the research. It defines the formal representation of the image data, the mathematical formulation of the geometric and pixel-level augmentations utilized, the precise derivations of the evaluation metrics (IoU, Dice, Precision, Recall), and the mathematical definitions of the applied loss functions.

\textbf{Chapter 4: Proposed Methodology} provides a deep, block-by-block architectural breakdown of the implemented systems. It details the dense skip pathways and ResNet50 encoder of the UNet++ models, the Atrous Spatial Pyramid Pooling mechanics of DeepLabV3+, the training pipeline algorithms, and the specific iterative algorithms employed for post-segmentation morphological skeletonization.

\textbf{Chapter 5: Experimental Design} outlines the rigorous experimental setup. It details the composition and statistical distribution of the 637-image dataset (500 train, 137 test), the specific hyperparameter configurations, the learning rate scheduling strategies, and the hardware/software ecosystem utilized to execute the experiments.

\textbf{Chapter 6: Results and Discussion} presents the empirical findings of the research. It includes extensive comparative tables of quantitative metrics, plots of the training and validation learning curves, and qualitative visual analyses of the predicted segmentation overlays. Crucially, it discusses the outputs of the skeletonization algorithm, analyzing the distinct vessel counts extracted from the test set.

Finally, \textbf{Chapter 7: Conclusion and Future Work} summarizes the core findings and contributions of the thesis, acknowledges current limitations (such as the handling of severely fragmented vessels), and proposes specific, actionable avenues for future research, including the transition to 3D volumetric segmentation and real-time clinical deployment.
